\chapter{2008年天津卷（文）}

\section{单选题}

\begin{problem}
设集合 \(U = \{x \in \N \mid 0 < x \leqslant 8\}\)，\(S = \{1, 2, 4, 5\}\)，\(T = \{3, 5, 7\}\)，则 \(S \cap (\complement_U T) =\)（\quad）

\choices
    {\(\{1,2,4\}\)}
    {\(\{1,2,3,4,5,7\}\)}
    {\(\{1,2\}\)}
    {\(\{1,2,4,5,6,8\}\)}
\end{problem}

\begin{answer}
A
\end{answer}

\begin{solution}
由全集的定义，\(\complement_U T=\{1,2,4,6,8\}\)，所以
\(S\cap(\complement_U T)=\{1,2,4,5\}\cap\{1,2,4,6,8\}=\{1,2,4\}\)，故选 A.
\end{solution}

\begin{problem}
设变量 \(x\)，\(y\) 满足约束条件 \(\begin{cases}
x - y \geqslant 0,\\
x + y \leqslant 1,\\
x + 2y \geqslant 1,
\end{cases}\) 则目标函数 \(z = 5x + y\) 的最大值（\quad）

\choices
    {\(2\)}
    {\(3\)}
    {\(4\)}
    {\(5\)}
\end{problem}

\begin{answer}
D
\end{answer}

\begin{solution}
三个约束条件所表示的边界直线分别为 \(y=x\)、\(x+y=1\) 和 \(x+2y=1\). 它们两两相交得到
\(\left(\frac{1}{2},\frac{1}{2}\right)\)、\(\left(\frac{1}{3},\frac{1}{3}\right)\) 和 \((1,0)\)，且这三个点构成可行域的顶点. 目标函数在三角形可行域上的最大值在顶点处取得. 代入得
\(z\left(\frac{1}{2},\frac{1}{2}\right)=3\)、\(z\left(\frac{1}{3},\frac{1}{3}\right)=2\)、\(z(1,0)=5\)，所以最大值为 \(5\)，故选 D.
\end{solution}

\begin{problem}
函数 \( y = 1 + \sqrt{x} \)（\(0 \leqslant x \leqslant 4\)）的反函数是（\quad）

\choices
    {\( y = (x - 1)^{2} \)（\(1 \leqslant x \leqslant 3\)）}
    {\( y = (x - 1)^{2} \)（\(0 \leqslant x \leqslant 4\)）}
    {\( y = x^{2} - 1 \)（\(1 \leqslant x \leqslant 3\)）}
    {\( y = x^{2} - 1 \)（\(0 \leqslant x \leqslant 4\)）}
\end{problem}

\begin{answer}
A
\end{answer}

\begin{solution}
由 \(0\leqslant x\leqslant4\)，可知 \(1\leqslant y=1+\sqrt{x}\leqslant3\)，且原函数在定义域上单调递增，因而存在反函数. 由
\(y-1=\sqrt{x}\) 得 \(x=(y-1)^2\)，交换 \(x\)，\(y\) 得反函数为 \(y=(x-1)^2\)，定义域为 \(1\leqslant x\leqslant3\)，故选 A.
\end{solution}

\begin{problem}
若等差数列 \(\{a_{n}\}\) 的前5项和 \(S_{5} = 25\)，且 \(a_{2} = 3\)，则 \(a_{7} =\)（\quad）

\choices
    {\(12\)}
    {\(13\)}
    {\(14\)}
    {\(15\)}
\end{problem}

\begin{answer}
B
\end{answer}

\begin{solution}
等差数列的前五项和为 \(S_5=5a_3\)，所以 \(a_3=5\). 由 \(a_2=3\) 得公差 \(d=a_3-a_2=2\)，从而 \(a_7=a_3+4d=5+8=13\)，故选 B.
\end{solution}

\begin{problem}
设 \(a\)，\(b\) 是两条直线，\(\alpha\)，\(\beta\) 是两个平面，则 \(a \perp b\) 的一个充分条件是（\quad）

\choices
    {\(a \perp \alpha\)，\(b \parallel \beta\)，\(\alpha \perp \beta\)}
    {\(a \perp \alpha\)，\(b \perp \beta\)，\(\alpha \parallel \beta\)}
    {\(a \subset \alpha\)，\(b \perp \beta\)，\(\alpha \parallel \beta\)}
    {\(a \subset \alpha\)，\(b \parallel \beta\)，\(\alpha \perp \beta\)}
\end{problem}

\begin{answer}
C
\end{answer}

\begin{solution}
由 \(\alpha\parallel\beta\)，且 \(b\perp\beta\)，可知 \(b\perp\alpha\). 又因 \(a\subset\alpha\)，所以 \(a\perp b\). 因此选项 C 给出了 \(a\perp b\) 的一个充分条件，故选 C.
\end{solution}

\begin{problem}
把函数 \( y = \sin x \)（\( x \in \R \)）的图象上所有的点向左平行移动 \( \frac{\pi}{3} \) 个单位长度，再把所得图象上所有点的横坐标缩短到原来的 \( \frac{1}{2} \) 倍（纵坐标不变），得到的图象所表示的函数是（\quad）

\choices
    {\( y = \sin \left(2x - \frac{\pi}{3}\right) \)（\(x \in \R\)）}
    {\(y = \sin \left(\frac{x}{2} + \frac{\pi}{6}\right)\)（\(x \in \R\)）}
    {\(y = \sin \left(2x + \frac{\pi}{3}\right)\)（\(x \in \R\)）}
    {\(y = \sin \left(2x + \frac{2\pi}{3}\right)\)（\(x \in \R\)）}
\end{problem}

\begin{answer}
C
\end{answer}

\begin{solution}
将 \(y=\sin x\) 的图象向左平移 \(\frac{\pi}{3}\) 后，所得函数为 \(y=\sin\left(x+\frac{\pi}{3}\right)\). 横坐标缩短到原来的 \(\frac{1}{2}\) 倍，等价于把自变量替换为 \(2x\)，所以所得函数为
\(y=\sin\left(2x+\frac{\pi}{3}\right)\)，故选 C.
\end{solution}

\begin{problem}
设椭圆 \(\frac{x^2}{m^2} + \frac{y^2}{n^2} = 1\)（\(m > 0\)，\(n > 0\)）的右焦点与抛物线 \(y^2 = 8x\) 的焦点相同，离心率为 \(\frac{1}{2}\)，则此椭圆的方程为（\quad）

\choices
    {\(\frac{x^2}{12} +\frac{y^2}{16} = 1\)}
    {\(\frac{x^2}{16} +\frac{y^2}{12} = 1\)}
    {\(\frac{x^2}{48} +\frac{y^2}{64} = 1\)}
    {\(\frac{x^2}{64} +\frac{y^2}{48} = 1\)}
\end{problem}

\begin{answer}
B
\end{answer}

\begin{solution}
抛物线 \(y^2=8x\) 可写为 \(y^2=4px\)，故 \(p=2\)，其焦点为 \((2,0)\). 因此椭圆的右焦距 \(c=2\). 由离心率 \(e=\frac{c}{m}=\frac{1}{2}\) 得 \(m=4\)，再由 \(c^2=m^2-n^2\) 得 \(n^2=16-4=12\). 所以椭圆方程为 \(\frac{x^2}{16}+\frac{y^2}{12}=1\)，故选 B.
\end{solution}

\begin{problem}
已知函数 \( f(x) = \begin{cases}
x + 2, & x \leqslant 0,\\
-x + 2, & x > 0,
\end{cases} \) 则不等式 \( f(x) \geqslant x^2 \) 的解集为（\quad）

\choices
    {\([-1, 1]\)}
    {\([-2, 2]\)}
    {\([-2, 1]\)}
    {\([-1, 2]\)}
\end{problem}

\begin{answer}
A
\end{answer}

\begin{solution}
当 \(x\leqslant0\) 时，\(f(x)\geqslant x^2\) 等价于
\(x+2\geqslant x^2\)，即 \((x-2)(x+1)\leqslant0\)，再结合 \(x\leqslant0\)，得 \(-1\leqslant x\leqslant0\).
当 \(x>0\) 时，\(f(x)\geqslant x^2\) 等价于 \(-x+2\geqslant x^2\)，即 \((x+2)(x-1)\leqslant0\)，再结合 \(x>0\)，得 \(0<x\leqslant1\). 合并得解集为 \([-1,1]\)，故选 A.
\end{solution}

\begin{problem}
设 \(a = \sin \frac{5\pi}{7}\)，\(b = \cos \frac{2\pi}{7}\)，\(c = \tan \frac{2\pi}{7}\)，则（\quad）

\choices
    {\(a < b < c\)}
    {\(a < c < b\)}
    {\(b < c < a\)}
    {\(b < a < c\)}
\end{problem}

\begin{answer}
D
\end{answer}

\begin{solution}
令 \(\theta=\frac{2\pi}{7}\)，则 \(a=\sin\frac{5\pi}{7}=\sin\frac{2\pi}{7}=\sin\theta\)，\(b=\cos\theta\)，\(c=\tan\theta\). 由于 \(\frac{\pi}{4}<\theta<\frac{\pi}{2}\)，有 \(\sin\theta>\cos\theta\)，又因 \(0<\cos\theta<1\)，有 \(\tan\theta=\frac{\sin\theta}{\cos\theta}>\sin\theta\). 所以 \(b<a<c\)，故选 D.
\end{solution}

\begin{problem}
设 \(a > 1\)，若对于任意的 \(x \in [a, 2a]\)，都有 \(y \in [a, a^2]\) 满足方程 \(\log_a x + \log_a y = 3\). 这时 \(a\) 的取值的集合为（\quad）

\choices
    {\(\{a\mid 1 < a \leqslant 2\}\)}
    {\(\{a\mid a \geqslant 2\}\)}
    {\(\{a\mid 2 \leqslant a \leqslant 3\}\)}
    {\(\{2, 3\}\)}
\end{problem}

\begin{answer}
B
\end{answer}

\begin{solution}
由 \(\log_a x+\log_a y=3\)，得 \(xy=a^3\)，所以对给定的 \(x\)，必须取 \(y=\frac{a^3}{x}\). 当 \(x\in[a,2a]\) 时，函数 \(\frac{a^3}{x}\) 递减，故其取值范围为
\(\left[\frac{a^2}{2},a^2\right]\). 要使任意这样的 \(y\) 都属于 \([a,a^2]\)，只需且必须有 \(\frac{a^2}{2}\geqslant a\). 因为 \(a>1\)，这等价于 \(a\geqslant2\). 故取值集合为 \(\{a\mid a\geqslant2\}\)，选 B.
\end{solution}

\section{填空题}

\begin{problem}
一个单位共有职工200人，其中不超过45岁的有120人，超过45岁的有80人. 为了调查职工的健康状况，用分层抽样的方法从全体职工中抽取一个容量为25的样本，应抽取超过45岁的职工\fillinblank{}人.
\end{problem}

\begin{answer}
10
\end{answer}

\begin{solution}
超过 \(45\) 岁的职工占全体职工的比例为 \(\frac{80}{200}=\frac{2}{5}\). 按分层抽样，样本中应抽取超过 \(45\) 岁的职工人数为 \(25\times\frac{2}{5}=10\).
\end{solution}

\begin{problem}
\(\left(x + \frac{2}{x}\right)^5\) 的二项展开式中 \(x^3\) 的系数为\fillinblank{}.（用数字作答）
\end{problem}

\begin{answer}
10
\end{answer}

\begin{solution}
二项展开式的通项为 \(\mathrm C_5^k x^{5-k}\left(\frac{2}{x}\right)^k=\mathrm C_5^k2^k x^{5-2k}\). 令 \(5-2k=3\)，得 \(k=1\)，所以 \(x^3\) 的系数为 \(\mathrm C_5^1\cdot2=10\).
\end{solution}

\begin{problem}
若一个球的体积为 \(4\sqrt{3}\pi\)，则它的表面积为\fillinblank{}.
\end{problem}

\begin{answer}
\(12\pi\)
\end{answer}

\begin{solution}
设球的半径为 \(r\). 由球的体积公式，\(\frac{4}{3}\pi r^3=4\sqrt{3}\pi\)，得 \(r^3=3\sqrt{3}=(\sqrt{3})^3\)，所以 \(r=\sqrt{3}\). 因此球的表面积为 \(4\pi r^2=4\pi\cdot3=12\pi\).
\end{solution}

\begin{problem}
已知平面向量 \(\bs{a} = (2, 4)\)，\(\bs{b} = (-1, 2)\)，若 \(\bs{c} = \bs{a} - (\bs{a} \cdot \bs{b})\bs{b}\). 则 \(|\bs{c}| = \)\fillinblank{}.
\end{problem}

\begin{answer}
\(8\sqrt{2}\)
\end{answer}

\begin{solution}
先求数量积：\(\bs{a}\cdot\bs{b}=2\times(-1)+4\times2=6\). 于是
\(\bs{c}=\bs{a}-6\bs{b}=(2,4)-6(-1,2)=(8,-8)\)，所以 \(|\bs{c}|=\sqrt{8^2+(-8)^2}=8\sqrt{2}\).
\end{solution}

\begin{problem}
已知圆 \(C\) 的圆心与点 \(P(-2,1)\) 关于直线 \(y = x + 1\) 对称. 直线 \(3x + 4y - 11 = 0\) 与圆 \(C\) 相交于 \(A\)，\(B\) 两点，且 \(|AB| = 6\)，则圆 \(C\) 的方程为\fillinblank{}.
\end{problem}

\begin{answer}
\(x^2 + (y + 1)^2 = 18\)
\end{answer}

\begin{solution}
设圆心为 \(O\). 将 \(P(-2,1)\) 关于直线 \(x-y+1=0\) 对称，利用点关于直线的对称公式，得 \(O=(0,-1)\). 直线 \(3x+4y-11=0\) 到圆心的距离为 \(d=\frac{|3\times0+4\times(-1)-11|}{\sqrt{3^2+4^2}}=3\). 弦 \(AB\) 的半长为 \(3\)，故在圆心、弦中点和弦端点构成的直角三角形中，圆半径满足 \(r^2=d^2+3^2=18\). 所以圆 \(C\) 的方程为 \(x^2+(y+1)^2=18\).
\end{solution}

\begin{problem}
有 4 张分别标有数字 1， 2， 3， 4 的红色卡片和 4 张分别标有数字 1， 2， 3， 4 的蓝色卡片，从这 8 张卡片中取出 4 张卡片排成一行. 如果取出的 4 张卡片所标的数字之和等于 10，则不同的排法共有\fillinblank{}种.（用数字作答）
\end{problem}

\begin{answer}
432
\end{answer}

\begin{solution}
先按数字统计排列类型. 4 个数字之和为 \(10\)，且同一数字最多出现两次（因为同一数字只有红、蓝两张卡片），可能的多重集合只有
\(\{1,1,4,4\}\)、\(\{1,2,3,4\}\)、\(\{1,3,3,3\}\)、\(\{2,2,2,4\}\)、\(\{2,2,3,3\}\). 其中含有数字出现三次的两类不能取出，剩下三类分别计数：\(\{1,2,3,4\}\) 有 \(4!\) 种数字排列，每个数字有 \(2\) 种颜色可选，共 \(4!\cdot2^4=384\) 种；\(\{1,1,4,4\}\) 和 \(\{2,2,3,3\}\) 各有 \(\frac{4!}{2!2!}=6\) 种数字排列，每个重复数字的两张卡片颜色必须一红一蓝，共 \(6\cdot2\cdot2=24\) 种，两个类型合计 \(48\) 种. 因此不同排法共有 \(384+48=432\) 种.
\end{solution}

\section{解答题}

\begin{problem}
已知函数 \( f(x) = 2\cos^{2}\omega x + 2\sin\omega x\cos \omega x + 1 \)（\(x \in \R\)，\(\omega > 0\)）的最小正周期是 \(\frac{\pi}{2}\).
\begin{enumerate}
\item 求 \( \omega \) 的值；
\item 求函数 \( f(x) \) 的最大值，并且求使 \( f(x) \) 取得最大值的 \( x \) 的集合.
\end{enumerate}
\end{problem}

\begin{solution}
\begin{enumerate}
\item 利用二倍角公式，得
\[
f(x)=2+\cos(2\omega x)+\sin(2\omega x)=2+\sqrt{2}\sin\left(2\omega x+\frac{\pi}{4}\right)
\]
因此 \(f(x)\) 的最小正周期为 \(\frac{2\pi}{2\omega}=\frac{\pi}{\omega}\). 由题意 \(\frac{\pi}{\omega}=\frac{\pi}{2}\)，所以 \(\omega=2\).
\item 由（1），\(f(x)=2+\sqrt{2}\sin\left(4x+\frac{\pi}{4}\right)\). 因为 \(-1\leqslant\sin\left(4x+\frac{\pi}{4}\right)\leqslant1\)，所以 \(f(x)\) 的最大值为 \(2+\sqrt{2}\). 等号成立当且仅当
\(4x+\frac{\pi}{4}=\frac{\pi}{2}+2k\pi\)，\(k\in\Z\)
即 \(x=\frac{\pi}{16}+\frac{k\pi}{2}\). 故使 \(f(x)\) 取得最大值的 \(x\) 的集合为 \(\left\{x\mathrel{\left|\right.}x=\frac{\pi}{16}+\frac{k\pi}{2},\ k\in\Z\right\}\).
\end{enumerate}
\end{solution}

\begin{problem}
甲，乙两个篮球运动员互不影响地在同一位置投球，命中率分别为 \(\frac{1}{2}\) 与 \(p\)，且乙投球 2 次均未命中的概率为 \(\frac{1}{16}\).
\begin{enumerate}
\item 求乙投球的命中率 \(p\)；
\item 求甲投球 2 次，至少命中 1 次的概率；
\item 若甲，乙两人各投球 2 次，求两人共命中 2 次的概率.
\end{enumerate}
\end{problem}

\begin{solution}
\begin{enumerate}
\item 设乙第 \(i\) 次投球命中为事件 \(B_i\)（\(i=1\)，\(2\)），则 \(P(B_i)=p\). 由两次投球相互独立，乙两次均未命中的概率为
\[
P(\overline{B_1}\cap\overline{B_2})=(1-p)^2=\frac{1}{16}
\]
因为 \(0\leqslant p\leqslant1\)，所以 \(1-p=\frac{1}{4}\)，从而 \(p=\frac{3}{4}\).
\item 设甲第 \(i\) 次投球命中为事件 \(A_i\)（\(i=1\)，\(2\)），则 \(P(A_i)=\frac{1}{2}\)，\(P(\overline{A_i})=\frac{1}{2}\). 甲投球 2 次至少命中 1 次的概率为
\[
1-P(\overline{A_1}\cap\overline{A_2})=1-\left(\frac{1}{2}\right)^2=\frac{3}{4}
\]
\item 甲、乙两人各投球 2 次时，共命中 2 次有三种互斥情形：甲命中 0 次且乙命中 2 次，甲命中 1 次且乙命中 1 次，甲命中 2 次且乙命中 0 次. 因此所求概率为
\[
\left(\frac{1}{2}\right)^2\left(\frac{3}{4}\right)^2+\mathrm C_2^1\frac{1}{2}\frac{1}{2}\cdot\mathrm C_2^1\frac{3}{4}\frac{1}{4}+\left(\frac{1}{2}\right)^2\left(\frac{1}{4}\right)^2=\frac{9}{64}+\frac{3}{16}+\frac{1}{64}=\frac{11}{32}
\]
\end{enumerate}
\end{solution}

\begin{problem}
如图，在四棱锥 \(P - ABCD\) 中，底面 \(ABCD\) 是矩形. 已知 \(AB = 3\)，\(AD = 2\)，\(PA = 2\)，\(PD = 2\sqrt{2}\)，\(\angle PAB = 60^\circ\).
\begin{enumerate}
\item 证明 \(AD \perp\) 平面 \(PAB\)；
\item 求异面直线 \(PC\) 与 \(AD\) 所成的角的大小；
\item 求二面角 \(P - BD - A\) 的大小.
\end{enumerate}

\bitmapfigure[width=0.36\linewidth]{img/2008/tianjin_liberal/q19_fig1.png}
\end{problem}

\begin{solution}
\begin{enumerate}
\item 在 \(\triangle PAD\) 中，\(PA^2+AD^2=2^2+2^2=8=PD^2\)，所以 \(AD\perp PA\). 又因底面 \(ABCD\) 是矩形，\(AD\perp AB\). 直线 \(PA\)，\(AB\) 是平面 \(PAB\) 内相交的两条直线，故 \(AD\perp\) 平面 \(PAB\).
\item 因为矩形 \(ABCD\) 中 \(BC\parallel AD\)，所以异面直线 \(PC\) 与 \(AD\) 所成的角等于直线 \(PC\) 与 \(BC\) 所成的锐角. 在 \(\triangle PAB\) 中，由余弦定理，
\[
PB^2=PA^2+AB^2-2PA\cdot AB\cos60^\circ=4+9-6=7
\]
因为 \(PB\) 是平面 \(PAB\) 内的直线，由（1）知 \(AD\perp PB\)，从而 \(PB\perp BC\). 因此 \(\triangle PBC\) 在 \(B\) 处为直角三角形，且 \(BC=AD=2\). 设所成的锐角为 \(\theta\)，则 \(\tan\theta=\frac{PB}{BC}=\frac{\sqrt{7}}{2}\)，所以 \(\theta=\arctan\frac{\sqrt{7}}{2}\).
\item 过点 \(P\) 作 \(PH\perp AB\)，垂足为 \(H\)；过点 \(H\) 作 \(HE\perp BD\)，垂足为 \(E\). 在直角三角形 \(PAH\) 中，
\(PH=PA\sin60^\circ=\sqrt{3}\)，\(AH=PA\cos60^\circ=1\).
所以 \(BH=AB-AH=2\)，且 \(BD=\sqrt{AB^2+AD^2}=\sqrt{13}\). 由于 \(BH\subset AB\) 且 \(AD\perp AB\)，点 \(D\) 到直线 \(BH\) 的距离为 \(AD\)，故
\[
S_{\triangle HBD}=\frac{1}{2}BH\cdot AD
\]
又因 \(HE\perp BD\)，所以
\[
S_{\triangle HBD}=\frac{1}{2}BD\cdot HE
\]
从而 \(HE=\frac{AD\cdot BH}{BD}=\frac{4}{\sqrt{13}}\).
由（1），\(PH\perp AD\)，又 \(PH\perp AB\)，所以 \(PH\perp\) 平面 \(ABCD\). 于是 \(PE\) 在平面 \(PBD\) 内且 \(PE\perp BD\)，\(HE\) 在平面 \(ABD\) 内且 \(HE\perp BD\)，故二面角 \(P-BD-A\) 的平面角为 \(\angle PEH\). 在直角三角形 \(PHE\) 中，
\[
\tan\angle PEH=\frac{PH}{HE}=\frac{\sqrt{3}}{4/\sqrt{13}}=\frac{\sqrt{39}}{4}
\]
所以二面角 \(P-BD-A\) 的大小为 \(\arctan\frac{\sqrt{39}}{4}\).
\end{enumerate}
\end{solution}

\begin{problem}
已知数列 \(\{a_{n}\}\) 中，\(a_{1} = 1\)，\(a_{2} = 2\)，且 \(a_{n+1} = (1 + q)a_{n} - qa_{n-1}\)（\(n \geqslant 2\)，\(q \neq 0\)）.
\begin{enumerate}
\item 设 \( b_{n} = a_{n+1} - a_{n} \)（\(n \in \N^{*}\)），证明 \( \{b_{n}\} \) 是等比数列；
\item 求数列 \(\{a_{n}\}\) 的通项公式；
\item 若 \(a_{3}\) 是 \(a_{6}\) 与 \(a_{9}\) 的等差中项，求 \(q\) 的值，并证明：对任意的 \(n \in \N^{*}\)，\(a_{n}\) 是 \(a_{n + 3}\) 与 \(a_{n + 6}\) 的等差中项.
\end{enumerate}
\end{problem}

\begin{solution}
\begin{enumerate}
\item 由递推式，
\(b_n=a_{n+1}-a_n=q(a_n-a_{n-1})=qb_{n-1}\)，\(n\geqslant2\).
又 \(b_1=a_2-a_1=1\)，所以 \(\{b_n\}\) 是首项为 \(1\)、公比为 \(q\) 的等比数列.
\item 由（1），当 \(n\geqslant2\) 时，\(a_n-a_{n-1}=q^{n-2}\)，从而
\[
a_n=1+\sum_{j=0}^{n-2}q^j=
\begin{cases}
1+\dfrac{1-q^{n-1}}{1-q},&q\ne1,\\
n,&q=1.
\end{cases}
\]
当 \(n=1\) 时上式也成立，因此这就是数列 \(\{a_n\}\) 的通项公式.
\item 若 \(q=1\)，则 \(a_n=n\)，此时 \(2a_3=6\ne15=a_6+a_9\)，所以 \(q\ne1\). 由通项公式，条件 \(2a_3=a_6+a_9\) 等价于
\[
a_3-a_6=a_9-a_3
\]
即
\[
\frac{q^2(q^3-1)}{1-q}=\frac{q^2(1-q^6)}{1-q}
\]
由于 \(q\ne0\) 且 \(q\ne1\)，可约去公因式，得 \(q^6+q^3-2=0\). 令 \(u=q^3\)，则 \((u-1)(u+2)=0\). 因 \(q\ne1\)，所以 \(q^3=-2\)，即 \(q=-\sqrt[3]{2}\).

对任意 \(n\in\N^*\)，由通项公式有
\(a_n-a_{n+3}=\frac{q^{n-1}(q^3-1)}{1-q}\)，\(a_{n+6}-a_n=\frac{q^{n-1}(1-q^6)}{1-q}\).
由 \(q^3=-2\) 得 \(q^3-1=-3=1-q^6\)，故 \(a_n-a_{n+3}=a_{n+6}-a_n\)，即 \(a_n\) 是 \(a_{n+3}\) 与 \(a_{n+6}\) 的等差中项.
\end{enumerate}
\end{solution}

\begin{problem}
设函数 \( f(x) = x^4 + ax^3 + 2x^2 + b \)（\(x \in \R\)），其中 \(a\)，\(b \in \R\).
\begin{enumerate}
\item 当 \(a = -\frac{10}{3}\) 时，讨论函数 \(f(x)\) 的单调性；
\item 若函数 \( f(x) \) 仅在 \( x = 0 \) 处有极值，求 \( a \) 的取值范围；
\item 若对于任意的 \(a \in [-2, 2]\)，不等式 \(f(x) \leqslant 1\) 在 \([-1, 1]\) 上恒成立，求 \(b\) 的取值范围.
\end{enumerate}
\end{problem}

\begin{solution}
\begin{enumerate}
\item 函数的导数为 \(f'(x)=4x^3+3ax^2+4x=x(4x^2+3ax+4)\). 当 \(a=-\frac{10}{3}\) 时，
\[
f'(x)=x(4x^2-10x+4)=2x(2x-1)(x-2)
\]
所以临界点为 \(0\)，\(\frac{1}{2}\)，\(2\). 逐区间判断导数符号可得：\(f'(x)<0\) 在 \((-\infty,0)\) 上成立，\(f'(x)>0\) 在 \(\left(0,\frac{1}{2}\right)\) 上成立，\(f'(x)<0\) 在 \(\left(\frac{1}{2},2\right)\) 上成立，\(f'(x)>0\) 在 \((2,+\infty)\) 上成立. 因此函数在 \((0,\frac{1}{2})\)，\((2,+\infty)\) 上单调递增，在 \((-\infty,0)\)，\(\left(\frac{1}{2},2\right)\) 上单调递减.
\item 记 \(g(x)=4x^2+3ax+4\)，则 \(f'(x)=xg(x)\). 要使 \(f(x)\) 除 \(x=0\) 外没有极值，二次函数 \(g(x)\) 不能有两个不同的实根；当它有一个二重根时，\(g(x)\) 不变号，该根也不是极值点. 因此必须且只需
\[
\Delta=(3a)^2-4\cdot4\cdot4=9a^2-64\leqslant0
\]
解得 \(-\frac{8}{3}\leqslant a\leqslant\frac{8}{3}\). 在此范围内，\(f'(x)\) 只在 \(x=0\) 处发生正负号改变，所以条件成立.
\item 对固定的 \(x\in[-1,1]\)，当 \(a\in[-2,2]\) 变化时，\(ax^3\) 的最大值为 \(2|x|^3\). 于是
\[
x^4+ax^3+2x^2+b\leqslant1
\]
对所有 \(a\in[-2,2]\) 和 \(x\in[-1,1]\) 恒成立，当且仅当
\[
b+x^4+2|x|^3+2x^2\leqslant1
\]
令 \(t=|x|\in[0,1]\)，则 \(t^4+2t^3+2t^2\) 在 \([0,1]\) 上递增，最大值为 \(1+2+2=5\). 所以必须有 \(b+5\leqslant1\)，即 \(b\leqslant-4\). 当 \(b\leqslant-4\) 时上述不等式对所有允许的 \(a\)，\(x\) 均成立，故取值范围为 \((-\infty,-4]\).
\end{enumerate}
\end{solution}

\begin{problem}
已知中心在原点的双曲线 \(C\) 的一个焦点是 \(F_{1}(-3,0)\)，一条渐近线的方程是 \(\sqrt{5} x - 2y = 0\).
\begin{enumerate}
\item 求双曲线 \(C\) 的方程；
\item 若以 \(k\)（\(k \neq 0\)）为斜率的直线 \(l\) 与双曲线 C 相交于两个不同的点 \(M\)，\(N\)，线段 \(MN\) 的垂直平分线与两坐标轴围成的三角形的面积为 \(\frac{81}{2}\)，求 \(k\) 的取值范围.
\end{enumerate}
\end{problem}

\begin{solution}
\begin{enumerate}
\item 设双曲线的标准方程为 \(\frac{x^2}{a^2}-\frac{y^2}{b^2}=1\)，其中 \(a>0\)，\(b>0\). 由焦点为 \((-3,0)\)，得焦半距 \(c=3\)，所以 \(a^2+b^2=c^2=9\). 又渐近线斜率为 \(\frac{b}{a}=\frac{\sqrt{5}}{2}\)，即 \(b^2=\frac{5}{4}a^2\). 联立得 \(a^2=4\)，\(b^2=5\)，故双曲线方程为 \(\frac{x^2}{4}-\frac{y^2}{5}=1\).
\item 设直线 \(l\) 的方程为 \(y=kx+m\)，其中 \(k\ne0\). 将其代入双曲线方程，得
\[
(5-4k^2)x^2-8kmx-4m^2-20=0
\]
直线与双曲线有两个不同交点，故二次项系数不为零且判别式大于零，即
\(5-4k^2\ne0\)，\(m^2+5-4k^2>0\).
记 \(A=5-4k^2\). 由韦达定理，弦 \(MN\) 的中点 \(Q(x_0,y_0)\) 满足
\(x_0=\frac{4km}{A}\)，\(y_0=kx_0+m=\frac{5m}{A}\).
因为 \(MN\) 的斜率为 \(k\)，所以其中垂线方程为
\[
y-\frac{5m}{A}=-\frac{1}{k}\left(x-\frac{4km}{A}\right)
\]
即 \(y=-\frac{x}{k}+\frac{9m}{A}\). 它与两坐标轴的截距分别为 \(\left(\frac{9km}{A},0\right)\) 和 \(\left(0,\frac{9m}{A}\right)\). 由三角形面积条件，
\[
\frac{1}{2}\left|\frac{9km}{A}\right|\left|\frac{9m}{A}\right|=\frac{81}{2}
\]
从而 \(m^2=\frac{A^2}{|k|}\). 代入判别式条件，得
\[
\frac{A^2}{|k|}+A>0
\]
令 \(t=|k|>0\)，则上式等价于
\[
(4t^2-5)(4t^2-t-5)>0
\]
由于 \(4t^2-t-5=(4t-5)(t+1)\)，且 \(t+1>0\)，所以 \(0<t<\frac{\sqrt{5}}{2}\) 或 \(t>\frac{5}{4}\). 因此
\[
k\in\left(-\infty,-\frac{5}{4}\right)\cup\left(-\frac{\sqrt{5}}{2},0\right)\cup\left(0,\frac{\sqrt{5}}{2}\right)\cup\left(\frac{5}{4},+\infty\right)
\]
\end{enumerate}
\end{solution}
